\documentclass[12pt, a4paper]{article}
\usepackage[utf8]{inputenc}
\usepackage[T1]{fontenc}
\usepackage[romanian]{babel}
\usepackage{geometry}
\usepackage{hyperref}
\usepackage{booktabs}
\usepackage{longtable}
\usepackage{array}
\usepackage{enumitem}
\usepackage{fancyhdr}
\usepackage{graphicx}
\usepackage{titlesec}

\geometry{
    a4paper,
    left=2.5cm,
    right=2.5cm,
    top=3cm,
    bottom=3cm
}

\hypersetup{
    colorlinks=true,
    linkcolor=black,
    urlcolor=blue
}

\pagestyle{fancy}
\fancyhf{}
\rhead{\small Software Design Document for DockerfileGen}
\rfoot{\thepage}
\lfoot{\small Universitatea de Vest Timi\c{s}oara}

\begin{document}

% -------------------------------------------------------
%  Pagina de titlu
% -------------------------------------------------------
\begin{titlepage}
    \centering
    \vspace*{1cm}
    {\scshape\LARGE Software Design Document\par}
    \vspace{0.5cm}
    {\large for\par}
    \vspace{0.5cm}
    {\Huge\bfseries DockerfileGen\par}
    \vspace{1.5cm}
    {\Large Version 1.2\par}
    \vspace{1cm}
    {\bfseries Realizat de:\par}
    Apostol Andrei, Stan Alexandru, Lunculescu Gheorghe,\\
    Bolea Alexandra, C\^{i}rlugea Alexandru\par
    \vspace{1cm}
    {\bfseries Universitatea de Vest Timi\c{s}oara\par}
    \vspace{0.5cm}
    {\large 12.06.2026\par}
\end{titlepage}

% -------------------------------------------------------
%  Cuprins
% -------------------------------------------------------
\tableofcontents
\newpage

% -------------------------------------------------------
%  1. Introducere
% -------------------------------------------------------
\section{Introducere}

\subsection{Scopul documentului}

Acest document define\c{s}te arhitectura tehnic\u{a}, modulele de cod, structura de comunicare \c{s}i
implementarea sistemului DockerGen. Ac\c{t}ioneaz\u{a} ca un manual tehnic pentru dezvoltatorii care
doresc s\u{a} \^{i}n\c{t}eleag\u{a}, s\u{a} men\c{t}in\u{a} sau s\u{a} extind\u{a} aplica\c{t}ia.

\subsection{Aria de acoperire}

Documentul acoper\u{a} implementarea tehnic\u{a} a celor trei componente (client, server \c{s}i admin),
modul de interogare a API-urilor externe (Repology / Homebrew) \c{s}i detalieaz\u{a} protocolul de
re\c{t}ea personalizat utilizat pentru transmiterea instruc\c{t}iunilor, gestiunea fi\c{s}ierelor \c{s}i
controlul erorilor.

\subsection{Referin\c{t}e}

Documentul Software Requirements Specification (SRS) --- versiunea 1.2.

% -------------------------------------------------------
%  2. Prezentare generala
% -------------------------------------------------------
\section{Prezentare general\u{a} a sistemului}

\subsection{Descrierea sistemului}

DockerGen este o solu\c{t}ie distribuit\u{a} care automatizeaz\u{a} crearea de fi\c{s}iere Dockerfile.
Sistemul este compus din trei entit\u{a}\c{t}i: un client de execu\c{t}ie pentru generarea fi\c{s}ierelor,
un server central de procesare \c{s}i un panou de administrare dedicat monitoriz\u{a}rii \^{i}n timp real
a st\u{a}rii serviciului \c{s}i a clien\c{t}ilor conecta\c{t}i. Serverul interogheaz\u{a} surse externe
de \^{i}ncredere pentru a valida pachetele solicitate \c{s}i asambleaz\u{a} un fi\c{s}ier final optimizat.

\subsection{Obiective de design}

Izolarea sarcinilor de re\c{t}ea prin fire de execu\c{t}ie POSIX (\texttt{pthread}), asigurarea
integrit\u{a}\c{t}ii datelor prin protocolul TCP \c{s}i modularizarea logicii de parsare.
Comunica\c{t}ia inter-thread se realizeaz\u{a} printr-o coad\u{a} FIFO implementat\u{a} cu un pipe
anonim, garantând scrieri atomice.

\subsection{Arhitectura}

Sistemul utilizeaz\u{a} o arhitectur\u{a} distribuit\u{a} Multi-Client Server. Aceasta include un flux
de date TCP pentru generarea Dockerfile-urilor \c{s}i un flux de control UDP pentru func\c{t}iile de
administrare.

% -------------------------------------------------------
%  3. Design arhitectura
% -------------------------------------------------------
\section{Design arhitectur\u{a}}

\subsection*{Component breakdown}

Arhitectura proiectului este definit\u{a} de trei componente critice:

\begin{itemize}
    \item \textbf{A. Componenta Client (\texttt{client.c}):} Gestioneaz\u{a} interfata CLI, proceseaz\u{a}
    argumentele utilizatorului (\texttt{--dep}, \texttt{--env}, \texttt{--copy}, \texttt{--upload},
    \texttt{--get}, \texttt{--list}, \texttt{--delete}) \c{s}i scrie fluxul de date recep\c{t}ionat pe disc.

    \item \textbf{B. Componenta Server (\texttt{server.c}):} Entitatea central\u{a} care utilizeaz\u{a}
    \texttt{poll()} pentru multiplexarea conexiunilor TCP \c{s}i fire de execu\c{t}ie POSIX
    (\texttt{pthread}) pentru paralelizarea sarcinilor: un thread de procesare a cererilor (alimentat
    printr-un pipe anonim), un thread UDP de administrare, un thread \texttt{inotify} \c{s}i un thread
    de consol\u{a} interactiv\u{a} (\texttt{stdin}).

    \item \textbf{C. Componenta Admin (\texttt{admin.c}):} Un panou vizual (TUI) bazat pe
    \texttt{ncurses} care permite administratorului s\u{a} interogheze serverul prin UDP pentru a
    ob\c{t}ine statistici \c{s}i a gestiona clien\c{t}ii, f\u{a}r\u{a} a ocupa un slot de conexiune TCP.
\end{itemize}

\subsection*{Tehnologii utilizate}

\begin{itemize}
    \item \textbf{Limbaj \c{s}i Build:} Limbajul C, gestiune pachete: Conan, build system: CMake.
    \item \textbf{Biblioteci externe:}
    \begin{itemize}
        \item \texttt{ncurses}: pentru interfata vizual\u{a} (TUI) a panoului de administrare.
        \item \texttt{libcurl}: pentru comunicatia HTTP/JSON cu API-urile web.
        \item \texttt{libconfig}: pentru parsarea fi\c{s}ierelor de configurare \texttt{.cfg}.
        \item \texttt{pthread}: pentru toate mecanismele de concuren\c{t}\u{a} (fire de execu\c{t}ie,
        mutex-uri).
    \end{itemize}
    \item \textbf{Protocoale \c{s}i Mecanisme:}
    \begin{itemize}
        \item \textbf{TCP (portul 8080):} pentru transferul fi\c{s}ierelor \^{i}ntre client \c{s}i server.
        \item \textbf{UDP (portul 8081):} pentru schimbul rapid de comenzi de control \^{i}ntre admin
        \c{s}i server.
        \item \textbf{Multiplexare \texttt{poll()}:} pentru gestionarea eficient\u{a} a descriptorilor
        de socket \c{s}i suport multi-client \^{i}n server (timeout 500~ms).
        \item \textbf{Pipe anonim:} coad\u{a} FIFO \^{i}ntre thread-ul principal de I/O \c{s}i
        thread-ul de procesare; elementul transportat este \texttt{WorkItem}
        (\texttt{sizeof(WorkItem) < PIPE\_BUF}), garantând atomicitate.
        \item \textbf{\texttt{inotify}:} monitorizare \^{i}n timp real a directorului
        \texttt{uploads/} (evenimente \texttt{IN\_CREATE}, \texttt{IN\_MODIFY},
        \texttt{IN\_DELETE}, \texttt{IN\_MOVED\_TO}, \texttt{IN\_MOVED\_FROM},
        \texttt{IN\_CLOSE\_WRITE}).
    \end{itemize}
\end{itemize}

\subsection*{Fluxul de date}

Thread-ul principal prime\c{s}te cererea TCP \^{i}n bucla \texttt{poll()} $\rightarrow$ pentru
cererile de tip Dockerfile/LIST/DOWNLOAD/DELETE scrie un \texttt{WorkItem} \^{i}n pipe-ul anonim
$\rightarrow$ thread-ul de procesare cite\c{s}te din pipe \c{s}i apeleaz\u{a}
\texttt{process\_request\_and\_send()} sau handlerul corespunz\u{a}tor $\rightarrow$ pentru fiecare
pachet se lanseaz\u{a} un thread dedicat (\texttt{dep\_thread\_fn}) care interogheaz\u{a} API-urile
externe \c{s}i scrie rezultatul direct \^{i}n socket-ul clientului $\rightarrow$ thread-ul de
procesare trimite markerul de finalizare \texttt{===EOF===}.

Cererile de tip \texttt{UPLOAD} sunt tratate inline \^{i}n thread-ul I/O (streaming de date binare
f\u{a}r\u{a} a stoca fi\c{s}ierul \^{i}n memorie).

\begin{figure}[h]
    \centering
    \fbox{\parbox{0.85\textwidth}{\centering\vspace{2cm}[Diagrama de context]\vspace{2cm}}}
    \caption{Diagrama de context}
    \label{fig:context}
\end{figure}

\begin{figure}[h]
    \centering
    \fbox{\parbox{0.85\textwidth}{\centering\vspace{2cm}[Diagrama de arhitectur\u{a}]\vspace{2cm}}}
    \caption{Diagrama de arhitectur\u{a}}
    \label{fig:arch}
\end{figure}

% -------------------------------------------------------
%  4. Design detaliat
% -------------------------------------------------------
\section{Design detaliat}

\subsection{Modulul client (\texttt{client.c})}

\subsubsection{Responsabilit\u{a}\c{t}i}

Modulul ac\c{t}ioneaz\u{a} ca interfa\c{t}\u{a} direct\u{a} cu utilizatorul. Prelucreaz\u{a} comenzile
introduse \^{i}n terminal, le formateaz\u{a} conform protocolului de re\c{t}ea stabilit \c{s}i expediaz\u{a}
cererea c\u{a}tre server. Ulterior, intercepteaz\u{a} fluxul de date primit \c{s}i \^{i}l scrie pe discul
local \^{i}ntr-un format valid de infrastructur\u{a}. Suport\u{a} opera\c{t}iunile de gestiune a
fi\c{s}ierelor (upload, download, list, delete) pe acela\c{s}i socket TCP.

\subsubsection{Interfe\c{t}e \c{s}i API-uri}

Intr\u{a}rile sunt reprezentate de fluxul de caractere introdus de la tastatur\u{a} de c\u{a}tre
utilizator \c{s}i de fluxul binar recep\c{t}ionat din re\c{t}ea prin socket-ul TCP. Ie\c{s}irile
constau \^{i}n mesajele de status afi\c{s}ate pe consol\u{a} \c{s}i fi\c{s}ierul fizic generat la
loca\c{t}ia specificat\u{a} (implicit \texttt{Dockerfile.gen}).

Portul de conectare este citit din variabila de mediu \texttt{SERVER\_PORT} definit\u{a} \^{i}n
fi\c{s}ierul \texttt{.env}; valoarea implicit\u{a} este \textbf{8080}.

\subsubsection{Gestionarea erorilor}

Aceast\u{a} component\u{a} este conceput\u{a} s\u{a} trateze erorile la nivel local pentru a nu
supr\^{a}nc\u{a}rca inutil serverul cu trafic invalid. Dac\u{a} utilizatorul introduce o comand\u{a}
goal\u{a} sau una care nu respect\u{a} sintaxa argumentelor, modulul filtreaz\u{a} cererea \c{s}i
afi\c{s}eaz\u{a} imediat un ghid de utilizare, f\u{a}r\u{a} a ini\c{t}ia o nou\u{a} tranzac\c{t}ie
TCP. La nivel de re\c{t}ea, \^{i}n cazul imposibilit\u{a}\c{t}ii cre\u{a}rii socket-ului sau al
refuzului de conectare, clientul raporteaz\u{a} situa\c{t}ia \c{s}i se \^{i}nchide controlat.

\subsubsection{Structuri de date}

Componenta folose\c{s}te preponderent buffere de memorie prealocate (tablouri de caractere) pentru
stocarea temporar\u{a} a comenzii ce urmeaz\u{a} a fi expediat\u{a} \c{s}i pentru acumularea
pachetelor TCP primite de la server. Dimensiunile relevante: \texttt{MAX\_BUFFER\_SIZE = 1024},
\texttt{MAX\_PATH\_SIZE = 256}.

\subsubsection{Algoritmi \c{s}i logic\u{a}}

Nucleul de procesare al textului utilizeaz\u{a} func\c{t}ia \texttt{strtok} pentru a sec\c{t}iona
comanda ini\c{t}ial\u{a} \^{i}n func\c{t}ie de delimitatori. Pe m\u{a}sur\u{a} ce extrage argumentele,
algoritmul le concateneaz\u{a} cu prefixele protocolului intern (\texttt{D:}, \texttt{E:},
\texttt{C:}). Pentru mecanismul de sincronizare a salv\u{a}rii, o bucl\u{a} de recep\c{t}ie
analizeaz\u{a} continuu datele intrate folosind \texttt{strstr} pentru a detecta markerul
\texttt{===EOF===}, calculând exact pozi\c{t}ia de decupare a \c{s}irului pentru a exclude markerul
din fi\c{s}ierul final.

Pentru upload, clientul cite\c{s}te dimensiunea fi\c{s}ierului cu \texttt{stat()}, construie\c{s}te
header-ul \texttt{UPLOAD:basename:size} \c{s}i streameaz\u{a} con\c{t}inutul \^{i}n buc\u{a}\c{t}i
de 4096~bytes.

\subsubsection{Managementul st\u{a}rii}

Modulul opereaz\u{a} f\u{a}r\u{a} a men\c{t}ine o stare complex\u{a} \^{i}ntre gener\u{a}ri diferite.
Odat\u{a} ce un document a fost scris \c{s}i markerul de final a fost recep\c{t}ionat cu succes,
starea memoriei temporare este resetat\u{a}, permi\c{t}ând ini\c{t}ierea unei noi sesiuni de
comand\u{a} \^{i}n aceea\c{s}i instan\c{t}\u{a} a programului.

% ---
\subsection{Modulul server (\texttt{server.c})}

\subsubsection{Responsabilit\u{a}\c{t}i}

Acest modul reprezint\u{a} creierul sistemului. Coordoneaz\u{a} conexiunile de re\c{t}ea de pe portul
de ascultare TCP, cite\c{s}te propriet\u{a}\c{t}ile implicite din fi\c{s}ierul de configurare
\texttt{demo.cfg} (prin \texttt{libconfig}), parseaz\u{a} comenzile primite \c{s}i orchestreaz\u{a}
thread-urile care execut\u{a} interog\u{a}rile HTTP. \^{I}n final, asambleaz\u{a} toate
componentele \c{s}i le trimite \^{i}napoi c\u{a}tre client.

Serverul porne\c{s}te patru thread-uri de servicii la ini\c{t}ializare:

\begin{itemize}
    \item \textbf{\texttt{admin\_udp\_thread}} --- ascult\u{a} pe portul UDP 8081 \c{s}i gestioneaz\u{a}
    sesiunea exclusiv\u{a} 1:1 a administratorului.
    \item \textbf{\texttt{processing\_thread\_fn}} --- cite\c{s}te \texttt{WorkItem}-uri din pipe-ul
    anonim \c{s}i apeleaz\u{a} logica de asamblare Dockerfile sau handlerii de fi\c{s}iere.
    \item \textbf{\texttt{inotify\_thread\_fn}} --- monitorizeaz\u{a} directorul \texttt{uploads/}
    \c{s}i \^{i}nregistreaz\u{a} \^{i}n log evenimentele de sistem de fi\c{s}iere.
    \item \textbf{\texttt{stdin\_thread\_fn}} --- accept\u{a} comenzi interactive de la operatorul
    local: \texttt{help}, \texttt{status}, \texttt{clients}, \texttt{shutdown}.
\end{itemize}

\subsubsection{Interfe\c{t}e \c{s}i API-uri}

Intr\u{a}rile sistemului includ cererile TCP structurate venite de la client, r\u{a}spunsurile
formatate JSON preluate de la API-urile externe Repology \c{s}i Homebrew, precum \c{s}i datele statice
din fi\c{s}ierele locale \texttt{demo.cfg} \c{s}i \texttt{.env}. Ie\c{s}irea este exclusiv fluxul
text al fi\c{s}ierului asamblat, scris secven\c{t}ial direct \^{i}n socket-ul de retur al clientului.

\subsubsection{Gestionarea erorilor}

Modulul server exceleaz\u{a} prin abordarea erorilor cu o filozofie de notificare \textit{in-band},
axat\u{a} pe toleran\c{t}\u{a}. \^{I}n loc s\u{a} opreasc\u{a} brusc asamblarea dac\u{a} un API web
returneaz\u{a} o eroare 404 (pachet inexistent) sau dac\u{a} interogarea dep\u{a}\c{s}e\c{s}te limita
de timeout impus\u{a} prin \texttt{CURLOPT\_TIMEOUT} de 8~secunde, thread-ul de validare captureaz\u{a}
elegant acest e\c{s}ec. El scrie \^{i}n socket un comentariu de avertizare
(ex.\ \texttt{\# ESUAT: [pachet] (Lipseste de pe sursele verificate)}), permi\c{t}ând procesului
principal s\u{a} continue generarea documentului pentru restul componentelor valide.

Erorile de tip adres\u{a} blocat\u{a} sunt evitate prin configurarea socket-ului cu op\c{t}iunea
\texttt{SO\_REUSEADDR}.

\subsubsection{Structuri de date}

\begin{itemize}
    \item \textbf{\texttt{WorkItem}}: structur\u{a} transportat\u{a} prin pipe-ul anonim; con\c{t}ine
    \texttt{client\_fd} (int), \texttt{data[1024]} (char) \c{s}i \texttt{data\_len} (ssize\_t).
    \texttt{sizeof(WorkItem) < PIPE\_BUF}, garantând atomicitate.
    \item \textbf{\texttt{CurlBuffer}}: buffer fix de 8192~bytes (\texttt{CURL\_BUFFER\_SIZE}) pentru
    stocarea r\u{a}spunsurilor JSON primite prin HTTP; callback-ul \texttt{curl\_write\_buffer}
    trunchiaz\u{a} inteligent dac\u{a} se dep\u{a}\c{s}e\c{s}te limita.
    \item \textbf{\texttt{ClientInfo}}: stocheaz\u{a} fd-ul TCP \c{s}i adresa IP a fiec\u{a}rui
    client; tablou static de \texttt{MAX\_CLIENTS = 64} intr\u{a}ri, protejat de
    \texttt{g\_clients\_mutex}.
    \item \textbf{Starea sesiunii admin}: \texttt{g\_current\_admin\_addr} (sockaddr\_in),
    \texttt{g\_admin\_last\_activity} (time\_t), protejate de \texttt{g\_admin\_mutex}.
    \item \textbf{\texttt{DepArgs}}: argumente pasate thread-ului per-pachet: numele pachetului
    \c{s}i fd-ul clientului.
\end{itemize}

\subsubsection{Algoritmi \c{s}i logic\u{a}}

Serverul implementeaz\u{a} multiplexarea I/O utilizând func\c{t}ia \texttt{poll()} cu timeout de
500~ms, permi\c{t}ând gestionarea simultan\u{a} a mai multor descriptori de socket TCP f\u{a}r\u{a}
a bloca execu\c{t}ia.

Pentru validarea dependen\c{t}elor, se utilizeaz\u{a} un model paralel bazat pe thread-uri: pentru
fiecare pachet \texttt{D:} se creeaz\u{a} un thread dedicat (\texttt{dep\_thread\_fn}) care
interogheaz\u{a} mai \^{i}nt\^{a}i Repology, iar \^{i}n caz de e\c{s}ec, Homebrew. \^{I}ntre
interog\u{a}ri consecutive Repology se introduce o pauz\u{a} de 1~secund\u{a}
(\texttt{REPOLOGY\_DELAY}) pentru evitarea rate-limiting-ului (HTTP 429).

Un algoritm critic implementat manual este func\c{t}ia \texttt{extract\_json\_string}, care
analizeaz\u{a} semantic buffer-ul pentru a identifica cheia \texttt{"version"} (Repology) sau
\texttt{"stable"} (Homebrew), eliminând necesitatea integr\u{a}rii unor parsere externe.

Mecanismul de sesiune exclusiv\u{a} al administratorului func\c{t}ioneaz\u{a} astfel: la primirea
unui pachet UDP, thread-ul admin verific\u{a} sub mutex dac\u{a} sesiunea a expirat
(\texttt{>~60~secunde}) sau dac\u{a} expeditorul coincide cu sesiunea activ\u{a}. Dac\u{a} da,
actualizeaz\u{a} \texttt{g\_admin\_last\_activity}; altfel, returneaz\u{a} mesaj de eroare
f\u{a}r\u{a} a procesa comanda.

\subsubsection{Managementul st\u{a}rii}

Serverul este proiectat s\u{a} fie independent de stare \^{i}ntre cererile diferi\c{t}ilor clien\c{t}i.
Starea global\u{a} (lista de clien\c{t}i, sesiunea admin, \texttt{g\_running}) este protejat\u{a}
exclusiv prin mutex-uri. Variabila \texttt{g\_running} este de tip \texttt{volatile sig\_atomic\_t}
\c{s}i este setat\u{a} pe 0 de handlerul de semnal (\texttt{SIGINT}/\texttt{SIGTERM}) sau de comanda
\texttt{shutdown} din consola local\u{a}.

% ---
\subsection{Modulul Admin (\texttt{admin.c})}

\subsubsection{Responsabilit\u{a}\c{t}i}

Interogarea serverului pentru ob\c{t}inerea raportului de status, listarea clien\c{t}ilor
conecta\c{t}i, eliminarea unui client (kick), deconectarea sesiunii de administrare \c{s}i
cur\u{a}\c{t}area directorului \texttt{uploads/}.

\subsubsection{Interfe\c{t}e \c{s}i API-uri}

Utilizeaz\u{a} \texttt{ncurses} (\texttt{initscr}, \texttt{cbreak}, \texttt{noecho},
\texttt{keypad}) pentru a desena meniuri interactive navigabile cu s\u{a}ge\c{t}ile tastaturii.
Comunicarea se face prin socket-uri UDP pe portul 8081 via func\c{t}ia \texttt{send\_command\_to\_server()}.

Meniul expune \textbf{7 op\c{t}iuni}:

\begin{center}
\begin{tabular}{lll}
\toprule
\textbf{Op\c{t}iune} & \textbf{Comand\u{a} UDP} & \textbf{Descriere} \\
\midrule
1 & \texttt{CMD:STATUS}  & Raport status server (port, clien\c{t}i, uptime) \\
2 & \texttt{CMD:CLIENTS} & Lista clien\c{t}ilor TCP conecta\c{t}i \\
3 & \texttt{CMD:KICK}    & Elimin\u{a} primul client TCP activ \\
4 & \texttt{CMD:LOGOUT}  & Elibereaz\u{a} sesiunea de admin \c{s}i \^{i}nchide panoul \\
5 & \texttt{CMD:VERSION} & Versiunea serverului \\
6 & \texttt{CMD:CLEAN}   & \c{S}terge toate fi\c{s}ierele din \texttt{uploads/} \\
7 & \texttt{CMD:PING}    & Test de conectivitate \\
\bottomrule
\end{tabular}
\end{center}

\subsubsection{Gestionarea erorilor}

Implementeaz\u{a} un mecanism de timeout de 2~secunde (\texttt{SO\_RCVTIMEO}) pentru a preveni
blocarea interfe\c{t}ei \^{i}n cazul \^{i}n care serverul nu r\u{a}spunde. \^{I}n caz de timeout,
buffer-ul de r\u{a}spuns este populat cu mesajul \textit{``Eroare: Serverul nu a r\u{a}spuns sau
este offline''}, afi\c{s}at direct \^{i}n panoul TUI.

\subsubsection{Structuri de date}

Utilizeaz\u{a} un buffer static \texttt{response\_buffer[2048]} pentru stocarea r\u{a}spunsului
curent de la server \c{s}i un tablou de pointeri la \c{s}iruri de caractere (\texttt{options[7]})
pentru op\c{t}iunile meniului. Starea nav\u{a}rii este men\c{t}inut\u{a} printr-o variabil\u{a}
\texttt{choice} (int) \c{s}i un flag \texttt{running} (int).

\subsubsection{Algoritmi \c{s}i logic\u{a}}

Bucla principal\u{a} redeseneaz\u{a} ecranul la fiecare apel \texttt{getch()}, aplicând atributul
\texttt{A\_REVERSE} pe op\c{t}iunea curent\u{a} pentru eviden\c{t}ierea selec\c{t}iei. Afi\c{s}area
r\u{a}spunsului se face caracter cu caracter, cu gestionarea manual\u{a} a \textit{word-wrap}-ului
la \texttt{MAX\_LINE\_WIDTH = 78} coloane \c{s}i interpretarea caracterelor \texttt{'\textbackslash n'}
pentru salt de linie.

% -------------------------------------------------------
%  5. Interfete externe
% -------------------------------------------------------
\section{Interfe\c{t}e externe}

\subsection{Interfata utilizator}

Sistemul nu dispune de o interfa\c{t}\u{a} grafic\u{a} tradi\c{t}ional\u{a}, ci expune o interfa\c{t}\u{a}
\^{i}n linie de comand\u{a} strict\u{a}, optimizat\u{a} pentru ingineri software \c{s}i procese
automatizate. Promptul interactiv \texttt{comanda:>} indic\u{a} momentul \^{i}n care aplica\c{t}ia
este preg\u{a}tit\u{a} s\u{a} primeasc\u{a} o nou\u{a} serie de argumente.

Argumentele acceptate de client sunt prezentate \^{i}n tabelul urm\u{a}tor:

\begin{center}
\begin{tabular}{ll}
\toprule
\textbf{Argument} & \textbf{Descriere} \\
\midrule
\texttt{--dep <pachet>}    & Adaug\u{a} o dependen\c{t}\u{a} de validat \c{s}i instalat \\
\texttt{--env <KEY=VAL>}   & Adaug\u{a} o variabil\u{a} de mediu \^{i}n container \\
\texttt{--copy <src dst>}  & Specific\u{a} un fi\c{s}ier de copiat \^{i}n container \\
\texttt{--out <fisier>}    & Suprascriere nume implicit al fi\c{s}ierului generat \\
\texttt{--upload <cale>}   & \^{I}ncarc\u{a} un fi\c{s}ier local pe server \\
\texttt{--get <nume>}      & Desc\u{a}rc\u{a} un fi\c{s}ier de pe server \\
\texttt{--list}            & Afi\c{s}eaz\u{a} fi\c{s}ierele din \texttt{uploads/} al serverului \\
\texttt{--delete <nume>}   & \c{S}terge un fi\c{s}ier de pe server \\
\texttt{exit}              & \^{I}nchide clientul \\
\bottomrule
\end{tabular}
\end{center}

Pe l\^{a}ng\u{a} CLI-ul clientului, sistemul expune o interfa\c{t}\u{a} semi-grafic\u{a} (TUI) pentru
modulul de administrare, oferind meniuri navigabile din s\u{a}ge\c{t}ile tastaturii \c{s}i delimit\u{a}ri
clare pe ecran pentru afi\c{s}area r\u{a}spunsurilor.

\subsection{API-uri externe}

Pentru a garanta acurate\c{t}ea versiunilor software incluse \^{i}n documentul generat, aplica\c{t}ia
depinde de interogarea \^{i}n timp real a dou\u{a} platforme web distincte.

\begin{itemize}
    \item \textbf{Repology API (v1) --- surs\u{a} primar\u{a}:} Serverul trimite un request GET la
    \texttt{https://repology.org/api/v1/project/<pachet>} \c{s}i extrage câmpul \texttt{"version"}
    din r\u{a}spunsul JSON. Dac\u{a} array-ul returnat este gol (\texttt{[]}) sau http code-ul nu
    este 200, se trece la fallback.

    \item \textbf{Homebrew API --- surs\u{a} secundar\u{a} (fallback):} Serverul trimite un request GET
    la \texttt{https://formulae.brew.sh/api/formula/<pachet>.json} \c{s}i extrage câmpul
    \texttt{"stable"} din sec\c{t}iunea \texttt{"versions"}.
\end{itemize}

Ambele interac\c{t}iuni sunt intermediate de \texttt{libcurl} cu header
\texttt{User-Agent: DockerGen/1.0} \c{s}i timeout de 8~secunde (\texttt{CURLOPT\_TIMEOUT}).
Buffer-ul de r\u{a}spuns este limitat la \texttt{CURL\_BUFFER\_SIZE = 8192~bytes} pentru a preveni
dep\u{a}\c{s}irile de memorie.

\subsection{Protocoale de re\c{t}ea \c{s}i comunicare}

Sistemul utilizeaz\u{a} un model de comunicare dual-protocol:

\begin{enumerate}
    \item \textbf{TCP (portul 8080):} Utilizat pentru transferul sigur \c{s}i ordonat al fluxului
    de date Dockerfile \^{i}ntre client \c{s}i server. Prefixele de protocol sunt:

    \begin{center}
    \begin{tabular}{lll}
    \toprule
    \textbf{Prefix} & \textbf{Tip} & \textbf{Descriere} \\
    \midrule
    \texttt{D:}                   & dependen\c{t}\u{a}        & Pachet de validat \c{s}i instalat \\
    \texttt{E:}                   & variabil\u{a} de mediu    & Pereche \texttt{cheie=valoare} \\
    \texttt{C:}                   & fi\c{s}ier de copiat       & Cale surs\u{a} + destina\c{t}ie \\
    \texttt{UPLOAD:basename:size} & upload                    & Header urmat de date binare \\
    \texttt{DOWNLOAD:basename}    & download                  & Cerere fi\c{s}ier de pe server \\
    \texttt{LIST}                 & listare                   & Cerere list\u{a} fi\c{s}iere \\
    \texttt{DELETE:basename}      & \c{s}tergere               & Cerere \c{s}tergere fi\c{s}ier \\
    \bottomrule
    \end{tabular}
    \end{center}

    Finalul oric\u{a}rui r\u{a}spuns de la server este marcat cu \texttt{===EOF===}.

    \item \textbf{UDP (portul 8081):} Utilizat pentru mesaje scurte de control \c{s}i monitorizare
    \^{i}ntre admin \c{s}i server: \texttt{CMD:STATUS}, \texttt{CMD:CLIENTS}, \texttt{CMD:KICK},
    \texttt{CMD:LOGOUT}, \texttt{CMD:VERSION}, \texttt{CMD:PING}, \texttt{CMD:CLEAN}.
\end{enumerate}

La nivel extern, comunicarea dintre procesele serverului \c{s}i platformele de validare se
realizeaz\u{a} exclusiv prin protocolul securizat HTTPS peste portul 443.

% -------------------------------------------------------
%  6. Cerinte non-functionale
% -------------------------------------------------------
\section{Cerin\c{t}e non-func\c{t}ionale relevante pentru design}

\subsection{Performan\c{t}\u{a}}

Bucla de acceptare a conexiunilor TCP utilizeaz\u{a} \texttt{poll()} cu timeout de 500~ms pentru a
r\u{a}m\^{a}ne receptiv\u{a} la semnalele de oprire. \^{I}ntre interog\u{a}rile Repology
consecutive se introduce o pauz\u{a} de 1~secund\u{a} (\texttt{REPOLOGY\_DELAY}) pentru evitarea
rate-limiting-ului (HTTP 429). Timeout-ul maxim per interogare extern\u{a} este de 8~secunde.

\subsection{Securitate}

Toate numele de fi\c{s}iere primite prin comenzile \texttt{UPLOAD}, \texttt{DOWNLOAD} \c{s}i
\texttt{DELETE} sunt validate pentru a preveni atacuri de tip path traversal: se resping secven\c{t}ele
\texttt{..}, \texttt{/} \c{s}i \texttt{\textbackslash}. Buffer-ul curl este limitat la
\texttt{CURL\_BUFFER\_SIZE = 8192~bytes} pentru a preveni dep\u{a}\c{s}irile de memorie.

\subsection{Configurabilitate}

Comportamentul de baz\u{a} al sistemului (imaginea Docker, maintainer, workdir) poate fi modificat
exclusiv prin alterarea fi\c{s}ierului \texttt{demo.cfg}, f\u{a}r\u{a} recompilare. Portul TCP de
ascultare este configurat prin variabila \texttt{SERVER\_PORT} din \texttt{.env}.

\subsection{Jurnalizare}

Sistemul de jurnalizare thread-safe (\texttt{server\_log}) scrie simultan pe \texttt{stdout} \c{s}i
\^{i}n fi\c{s}ierul \texttt{server.log} cu trei niveluri de severitate: \texttt{LOG\_INFO},
\texttt{LOG\_WARN}, \texttt{LOG\_ERROR}. Accesul la fi\c{s}ierul de log \c{s}i la
\texttt{stdout} este protejat de \texttt{g\_log\_mutex}.

\end{document}
